\documentclass[11pt,a4paper]{article}
\usepackage[utf8]{inputenc}
\usepackage[T1]{fontenc}
\usepackage[spanish,es-minimal,es-noshorthands]{babel}
\AtBeginDocument{\renewcommand{\sin}{\operatorname{sen}}\renewcommand{\arcsin}{\operatorname{arcsen}}}
\usepackage{lmodern}
\usepackage[sfdefault,lining]{FiraSans}
\usepackage[a4paper,top=2.3cm,bottom=2.2cm,left=1.8cm,right=1.8cm,headheight=1.3cm,headsep=0.4cm,footskip=1cm]{geometry}
\usepackage{amsmath,amssymb}
\allowdisplaybreaks
\usepackage{icomma}
\usepackage{graphicx}
\usepackage[table]{xcolor}
\usepackage{tikz}
\usepackage{fontawesome5}
\usepackage{ccicons}
\usepackage{fancyhdr}
\usepackage{lastpage}
\usepackage{tasks}
\usepackage{needspace}
\usepackage[most]{tcolorbox}
\usepackage{array}
\usepackage[hidelinks]{hyperref}
\AtBeginDocument{\hypersetup{pdftitle={Ficha de ejercicios de la unidad 4: Problemas métricos. 2º Bach.},pdfsubject={Matemáticas II 2º Bach.},pdfkeywords={ángulos, distancias, perpendicular común, proyecciones, simétricos, áreas, volúmenes, plano mediador, planos bisectores, esfera, PAU}}}

\definecolor{navy}{HTML}{1B2748}
\definecolor{accent}{HTML}{E8603A}
\definecolor{teal}{HTML}{1F9E90}
\definecolor{soft}{HTML}{EEF1F8}
\definecolor{muted}{HTML}{5D6685}

\setlength{\parindent}{0pt}
\setlength{\parskip}{3pt}
\settasks{label=\alph*), label-format=\bfseries\color{navy}, label-width=1.3em, item-indent=1.8em, column-sep=0.6em, after-item-skip=5pt, before-skip=3pt, after-skip=4pt}

% ---------- capçalera i peu ----------
\newcommand{\llicencia}{\href{https://creativecommons.org/licenses/by-nc-sa/4.0/deed.es}{\ccbyncsa}\ \ Licencia Creative Commons Atribución-NoComercial-CompartirIgual 4.0 Internacional (CC BY-NC-SA 4.0)}
\pagestyle{fancy}
\fancyhf{}
\fancyhead[L]{\raisebox{-0.35\height}{\includegraphics[height=0.95cm]{logo-gva.png}}}
\fancyhead[C]{\small\color{muted}Ficha de ejercicios · Unidad 4: Problemas métricos · 2\textsuperscript{o} Bach.}
\fancyhead[R]{\raisebox{-0.35\height}{\includegraphics[height=0.95cm]{logo-vermellar.png}}}
\fancyfoot[L]{\scriptsize\color{muted}\llicencia}
\fancyfoot[R]{\small\color{muted}\thepage}
\renewcommand{\headrulewidth}{0.4pt}
\renewcommand{\headrule}{\color{muted!40}\hrule height \headrulewidth}
\fancypagestyle{primera}{\fancyhf{}\renewcommand{\headrulewidth}{0pt}
  \fancyfoot[L]{\scriptsize\color{muted}\llicencia}
  \fancyfoot[R]{\small\color{muted}\thepage}}

% ---------- nivells i exercicis ----------
\newcommand{\estrelles}[1]{{\color{accent}\ifcase#1\or\faStar\faStar[regular]\faStar[regular]\or\faStar\faStar\faStar[regular]\or\faStar\faStar\faStar\fi}}
\newcounter{exer}
\newcommand{\ex}[1]{\par\needspace{4\baselineskip}\medskip\refstepcounter{exer}\noindent{\color{navy}\bfseries\theexer.}\ \ }

\newcommand{\seccio}[2]{\par\needspace{10\baselineskip}\bigskip
  \begin{tcolorbox}[colback=navy,colframe=navy,coltext=white,arc=2mm,boxrule=0pt,left=3mm,right=3mm,top=1.2mm,bottom=1.2mm]
  {\large\bfseries #1\quad #2}
  \end{tcolorbox}\vspace{-2pt}}


\setlength{\emergencystretch}{3em}  % les fórmules llargues (equacions de rectes, ternes) passen a la línia següent

% quadre de recordatori (només a la fitxa d'exercicis)
\newcommand{\recorda}[1]{\par\smallskip\begin{tcolorbox}[colback=accent!7,colframe=accent!55,arc=2mm,boxrule=0.5pt,left=2.5mm,right=2.5mm,top=1mm,bottom=1mm]\small{\color{accent}\faLightbulb[regular]}\ \textbf{Recuerda:}\ #1\end{tcolorbox}}
% etiqueta dels exercicis de la PAU d'altres comunitats
\newcommand{\pau}[1]{{\color{teal}\small\bfseries\faUniversity\ PAU · #1.}\ }
\newcommand{\res}[1]{#1}\newcommand{\rest}[1]{#1}

% =====================================================================
% Capçalera FentMates de les fitxes: un color per a cada curs.
% Els colors són els de la pàgina inicial (LEVELS a assets/js/materials.js);
% el logo ve de _eines/logo/pdf/ (python3 _eines/logo/genera.py).
%
% Ús, al final del preàmbul de la unitat:
%   % =====================================================================
% Capçalera FentMates de les fitxes: un color per a cada curs.
% Els colors són els de la pàgina inicial (LEVELS a assets/js/materials.js);
% el logo ve de _eines/logo/pdf/ (python3 _eines/logo/genera.py).
%
% Ús, al final del preàmbul de la unitat:
%   % =====================================================================
% Capçalera FentMates de les fitxes: un color per a cada curs.
% Els colors són els de la pàgina inicial (LEVELS a assets/js/materials.js);
% el logo ve de _eines/logo/pdf/ (python3 _eines/logo/genera.py).
%
% Ús, al final del preàmbul de la unitat:
%   \input{../comu/fentmates.tex}
%   \fmcurs{1eso}
% i, al començament del document:
%   \capcalera{Matemàtiques}{1r ESO}{Fitxa d'exercicis · Unitat 2}{Divisibilitat}
% Opcional abans de \input: \def\fmneed{14} (espai mínim abans de cada secció).
% =====================================================================
\makeatletter
\def\fm@logodir{../../logo/pdf}
\providecommand{\fmneed}{10}
\def\fm@cursnom{}

% clau del curs, color i color del text damunt del color
\newcommand{\fm@curs}[3]{\@namedef{fm@color@#1}{#2}\@namedef{fm@text@#1}{#3}}
\fm@curs{1eso}{E6B84A}{navy}
\fm@curs{2eso}{3B82F6}{white}
\fm@curs{3eso}{FF7A50}{white}
\fm@curs{4eso-a}{A78BFA}{white}
\fm@curs{4eso-b}{6D28D9}{white}
\fm@curs{1batx-mi}{9B1C1C}{white}
\fm@curs{1batx-ccss}{EF6B6B}{white}
\fm@curs{2batx-mii}{1E6B47}{white}
\fm@curs{2batx-ccss}{2BB3A3}{white}

\newcommand{\fmcurs}[1]{%
  \@ifundefined{fm@color@#1}{\PackageError{fentmates}{Curs desconegut: #1}{}}{}%
  \edef\fm@tmp{\noexpand\definecolor{curs}{HTML}{\@nameuse{fm@color@#1}}%
               \noexpand\colorlet{curstext}{\@nameuse{fm@text@#1}}}\fm@tmp
  \colorlet{cursfosc}{curs!72!black}%
  \def\fm@logo{\fm@logodir/vertical-#1.pdf}%
  \def\fm@marca{\fm@logodir/marca-#1.pdf}}

% ---------- capçalera de la primera pàgina ----------
% #1 matèria (petit), #2 curs, #3 tipus de document i unitat, #4 títol
\newcommand{\capcalera}[4]{%
  \gdef\fm@cursnom{#2}%
  \par\noindent
  \begin{tcolorbox}[colback=curs,colframe=curs,boxrule=0pt,arc=2.2mm,width=3.5cm,height=2.15cm,
      valign=center,left=3mm,right=2mm,top=1mm,bottom=1mm,nobeforeafter,box align=center]
    \color{curstext}{\footnotesize #1}\par\vspace{1pt}{\fontsize{19}{21}\selectfont\bfseries #2\par}%
  \end{tcolorbox}\hspace{5mm}%
  \begin{minipage}[c]{\dimexpr\linewidth-3.5cm-5mm-2.9cm\relax}
    {\color{cursfosc}\small\bfseries\MakeUppercase{#3}}\par\vspace{2pt}
    {\color{navy}\fontsize{20}{23}\selectfont\bfseries #4\par}
  \end{minipage}\hfill
  \begin{minipage}[c]{2.6cm}\raggedleft\includegraphics[width=2.6cm]{\fm@logo}\end{minipage}\par
  \vspace{2pt}}

% ---------- barres de secció del color del curs ----------
\renewcommand{\seccio}[2]{\par\needspace{\fmneed\baselineskip}\bigskip
  \begin{tcolorbox}[colback=curs,colframe=curs,coltext=curstext,arc=2mm,boxrule=0pt,left=3mm,right=3mm,top=1.2mm,bottom=1.2mm]
  {\large\bfseries #1\quad #2}
  \end{tcolorbox}\vspace{-2pt}}

% ---------- estrelles de dificultat del color del curs (a Batxillerat no se'n posen) ----------
\@ifundefined{estrelles}{}{%
  \renewcommand{\estrelles}[1]{{\color{curs}\ifcase#1\or\faStar\faStar[regular]\faStar[regular]\or\faStar\faStar\faStar[regular]\or\faStar\faStar\faStar\fi}}}

% ---------- capçalera i peu de les pàgines ----------
% primera pàgina: logos institucionals al peu
\fancypagestyle{primera}{\fancyhf{}\renewcommand{\headrulewidth}{0pt}
  \fancyfoot[L]{\raisebox{-0.32\height}{\includegraphics[height=0.78cm]{logo-gva.png}}\hspace{5mm}%
    \raisebox{-0.32\height}{\includegraphics[height=0.78cm]{logo-vermellar.png}}}
  \fancyfoot[R]{\scriptsize\color{muted}\href{https://creativecommons.org/licenses/by-nc-sa/4.0/deed.ca}{\ccbyncsa}\ \ CC BY-NC-SA 4.0 · fentmates.org\qquad\small\thepage}}
% la resta: el símbol a l'esquerra i el curs a la dreta, amb el títol al centre com abans
\fancyhead[L]{\raisebox{-0.3\height}{\includegraphics[height=0.62cm]{\fm@marca}}}
\fancyhead[R]{\small\bfseries\color{cursfosc}\fm@cursnom}
\renewcommand{\headrule}{\color{curs!45}\hrule height \headrulewidth}
\makeatother

%   \fmcurs{1eso}
% i, al començament del document:
%   \capcalera{Matemàtiques}{1r ESO}{Fitxa d'exercicis · Unitat 2}{Divisibilitat}
% Opcional abans de \input: \def\fmneed{14} (espai mínim abans de cada secció).
% =====================================================================
\makeatletter
\def\fm@logodir{../../logo/pdf}
\providecommand{\fmneed}{10}
\def\fm@cursnom{}

% clau del curs, color i color del text damunt del color
\newcommand{\fm@curs}[3]{\@namedef{fm@color@#1}{#2}\@namedef{fm@text@#1}{#3}}
\fm@curs{1eso}{E6B84A}{navy}
\fm@curs{2eso}{3B82F6}{white}
\fm@curs{3eso}{FF7A50}{white}
\fm@curs{4eso-a}{A78BFA}{white}
\fm@curs{4eso-b}{6D28D9}{white}
\fm@curs{1batx-mi}{9B1C1C}{white}
\fm@curs{1batx-ccss}{EF6B6B}{white}
\fm@curs{2batx-mii}{1E6B47}{white}
\fm@curs{2batx-ccss}{2BB3A3}{white}

\newcommand{\fmcurs}[1]{%
  \@ifundefined{fm@color@#1}{\PackageError{fentmates}{Curs desconegut: #1}{}}{}%
  \edef\fm@tmp{\noexpand\definecolor{curs}{HTML}{\@nameuse{fm@color@#1}}%
               \noexpand\colorlet{curstext}{\@nameuse{fm@text@#1}}}\fm@tmp
  \colorlet{cursfosc}{curs!72!black}%
  \def\fm@logo{\fm@logodir/vertical-#1.pdf}%
  \def\fm@marca{\fm@logodir/marca-#1.pdf}}

% ---------- capçalera de la primera pàgina ----------
% #1 matèria (petit), #2 curs, #3 tipus de document i unitat, #4 títol
\newcommand{\capcalera}[4]{%
  \gdef\fm@cursnom{#2}%
  \par\noindent
  \begin{tcolorbox}[colback=curs,colframe=curs,boxrule=0pt,arc=2.2mm,width=3.5cm,height=2.15cm,
      valign=center,left=3mm,right=2mm,top=1mm,bottom=1mm,nobeforeafter,box align=center]
    \color{curstext}{\footnotesize #1}\par\vspace{1pt}{\fontsize{19}{21}\selectfont\bfseries #2\par}%
  \end{tcolorbox}\hspace{5mm}%
  \begin{minipage}[c]{\dimexpr\linewidth-3.5cm-5mm-2.9cm\relax}
    {\color{cursfosc}\small\bfseries\MakeUppercase{#3}}\par\vspace{2pt}
    {\color{navy}\fontsize{20}{23}\selectfont\bfseries #4\par}
  \end{minipage}\hfill
  \begin{minipage}[c]{2.6cm}\raggedleft\includegraphics[width=2.6cm]{\fm@logo}\end{minipage}\par
  \vspace{2pt}}

% ---------- barres de secció del color del curs ----------
\renewcommand{\seccio}[2]{\par\needspace{\fmneed\baselineskip}\bigskip
  \begin{tcolorbox}[colback=curs,colframe=curs,coltext=curstext,arc=2mm,boxrule=0pt,left=3mm,right=3mm,top=1.2mm,bottom=1.2mm]
  {\large\bfseries #1\quad #2}
  \end{tcolorbox}\vspace{-2pt}}

% ---------- estrelles de dificultat del color del curs (a Batxillerat no se'n posen) ----------
\@ifundefined{estrelles}{}{%
  \renewcommand{\estrelles}[1]{{\color{curs}\ifcase#1\or\faStar\faStar[regular]\faStar[regular]\or\faStar\faStar\faStar[regular]\or\faStar\faStar\faStar\fi}}}

% ---------- capçalera i peu de les pàgines ----------
% primera pàgina: logos institucionals al peu
\fancypagestyle{primera}{\fancyhf{}\renewcommand{\headrulewidth}{0pt}
  \fancyfoot[L]{\raisebox{-0.32\height}{\includegraphics[height=0.78cm]{logo-gva.png}}\hspace{5mm}%
    \raisebox{-0.32\height}{\includegraphics[height=0.78cm]{logo-vermellar.png}}}
  \fancyfoot[R]{\scriptsize\color{muted}\href{https://creativecommons.org/licenses/by-nc-sa/4.0/deed.ca}{\ccbyncsa}\ \ CC BY-NC-SA 4.0 · fentmates.org\qquad\small\thepage}}
% la resta: el símbol a l'esquerra i el curs a la dreta, amb el títol al centre com abans
\fancyhead[L]{\raisebox{-0.3\height}{\includegraphics[height=0.62cm]{\fm@marca}}}
\fancyhead[R]{\small\bfseries\color{cursfosc}\fm@cursnom}
\renewcommand{\headrule}{\color{curs!45}\hrule height \headrulewidth}
\makeatother

%   \fmcurs{1eso}
% i, al començament del document:
%   \capcalera{Matemàtiques}{1r ESO}{Fitxa d'exercicis · Unitat 2}{Divisibilitat}
% Opcional abans de \input: \def\fmneed{14} (espai mínim abans de cada secció).
% =====================================================================
\makeatletter
\def\fm@logodir{../../logo/pdf}
\providecommand{\fmneed}{10}
\def\fm@cursnom{}

% clau del curs, color i color del text damunt del color
\newcommand{\fm@curs}[3]{\@namedef{fm@color@#1}{#2}\@namedef{fm@text@#1}{#3}}
\fm@curs{1eso}{E6B84A}{navy}
\fm@curs{2eso}{3B82F6}{white}
\fm@curs{3eso}{FF7A50}{white}
\fm@curs{4eso-a}{A78BFA}{white}
\fm@curs{4eso-b}{6D28D9}{white}
\fm@curs{1batx-mi}{9B1C1C}{white}
\fm@curs{1batx-ccss}{EF6B6B}{white}
\fm@curs{2batx-mii}{1E6B47}{white}
\fm@curs{2batx-ccss}{2BB3A3}{white}

\newcommand{\fmcurs}[1]{%
  \@ifundefined{fm@color@#1}{\PackageError{fentmates}{Curs desconegut: #1}{}}{}%
  \edef\fm@tmp{\noexpand\definecolor{curs}{HTML}{\@nameuse{fm@color@#1}}%
               \noexpand\colorlet{curstext}{\@nameuse{fm@text@#1}}}\fm@tmp
  \colorlet{cursfosc}{curs!72!black}%
  \def\fm@logo{\fm@logodir/vertical-#1.pdf}%
  \def\fm@marca{\fm@logodir/marca-#1.pdf}}

% ---------- capçalera de la primera pàgina ----------
% #1 matèria (petit), #2 curs, #3 tipus de document i unitat, #4 títol
\newcommand{\capcalera}[4]{%
  \gdef\fm@cursnom{#2}%
  \par\noindent
  \begin{tcolorbox}[colback=curs,colframe=curs,boxrule=0pt,arc=2.2mm,width=3.5cm,height=2.15cm,
      valign=center,left=3mm,right=2mm,top=1mm,bottom=1mm,nobeforeafter,box align=center]
    \color{curstext}{\footnotesize #1}\par\vspace{1pt}{\fontsize{19}{21}\selectfont\bfseries #2\par}%
  \end{tcolorbox}\hspace{5mm}%
  \begin{minipage}[c]{\dimexpr\linewidth-3.5cm-5mm-2.9cm\relax}
    {\color{cursfosc}\small\bfseries\MakeUppercase{#3}}\par\vspace{2pt}
    {\color{navy}\fontsize{20}{23}\selectfont\bfseries #4\par}
  \end{minipage}\hfill
  \begin{minipage}[c]{2.6cm}\raggedleft\includegraphics[width=2.6cm]{\fm@logo}\end{minipage}\par
  \vspace{2pt}}

% ---------- barres de secció del color del curs ----------
\renewcommand{\seccio}[2]{\par\needspace{\fmneed\baselineskip}\bigskip
  \begin{tcolorbox}[colback=curs,colframe=curs,coltext=curstext,arc=2mm,boxrule=0pt,left=3mm,right=3mm,top=1.2mm,bottom=1.2mm]
  {\large\bfseries #1\quad #2}
  \end{tcolorbox}\vspace{-2pt}}

% ---------- estrelles de dificultat del color del curs (a Batxillerat no se'n posen) ----------
\@ifundefined{estrelles}{}{%
  \renewcommand{\estrelles}[1]{{\color{curs}\ifcase#1\or\faStar\faStar[regular]\faStar[regular]\or\faStar\faStar\faStar[regular]\or\faStar\faStar\faStar\fi}}}

% ---------- capçalera i peu de les pàgines ----------
% primera pàgina: logos institucionals al peu
\fancypagestyle{primera}{\fancyhf{}\renewcommand{\headrulewidth}{0pt}
  \fancyfoot[L]{\raisebox{-0.32\height}{\includegraphics[height=0.78cm]{logo-gva.png}}\hspace{5mm}%
    \raisebox{-0.32\height}{\includegraphics[height=0.78cm]{logo-vermellar.png}}}
  \fancyfoot[R]{\scriptsize\color{muted}\href{https://creativecommons.org/licenses/by-nc-sa/4.0/deed.ca}{\ccbyncsa}\ \ CC BY-NC-SA 4.0 · fentmates.org\qquad\small\thepage}}
% la resta: el símbol a l'esquerra i el curs a la dreta, amb el títol al centre com abans
\fancyhead[L]{\raisebox{-0.3\height}{\includegraphics[height=0.62cm]{\fm@marca}}}
\fancyhead[R]{\small\bfseries\color{cursfosc}\fm@cursnom}
\renewcommand{\headrule}{\color{curs!45}\hrule height \headrulewidth}
\makeatother

\fmcurs{2batx-mii}

\begin{document}
\thispagestyle{primera}

\capcalera{Matemáticas II}{2\textsuperscript{o} Bach.}{Ficha de ejercicios · Unidad 4}{\raggedright Problemas métricos}
\vspace{10pt}
Nombre y apellidos: \rule{6.6cm}{0.4pt}\hfill Grupo: \rule{1.6cm}{0.4pt}\hfill Fecha: \rule{2.2cm}{0.4pt}
\vspace{8pt}

\begin{tcolorbox}[colback=soft,colframe=soft,arc=2mm,left=3mm,right=3mm,top=1.5mm,bottom=1.5mm]
\small Resuelve los ejercicios en el cuaderno justificando cada paso. Haz un esbozo cuando te ayude a ver la situación y comprueba siempre los resultados (que la proyección está en el plano o en la recta, que el punto medio del segmento $PP'$ es el punto $M$\ldots). Los ejercicios marcados con {\color{teal}\bfseries\faUniversity\ PAU} son de las pruebas de acceso a la universidad de otras comunidades autónomas (2024--2026).
\end{tcolorbox}

% =====================================================
\seccio{1}{Ángulos entre rectas y planos}

\ex{1} Calcula el ángulo que forman las rectas:\begin{tasks}(1)
  \task \mbox{$r:\dfrac{x-1}{2}=\dfrac{y+1}{1}=\dfrac{z}{-2}$} \ y \ $s:(x,\,y,\,z)=(0,\,1,\,3)+\lambda\,(4,\,0,\,3)$
  \task $r:\left\{\begin{aligned}x+y&=1\\y-z&=2\end{aligned}\right.$ \ y \ $s:x=y=z$
\end{tasks}\recorda{\mbox{$\cos\alpha=\dfrac{|\vec d_r\cdot\vec d_s|}{|\vec d_r|\,|\vec d_s|}$}, siempre entre $0^\circ$ y $90^\circ$.}

\ex{1} Calcula el ángulo que forma la recta \mbox{$r:\dfrac{x-2}{1}=\dfrac{y}{-1}=\dfrac{z+1}{2}$} con cada uno de los planos siguientes:\begin{tasks}(2)
  \task $\pi_1:x+y+z=4$
  \task $\pi_2:2x-2y+4z=1$
\end{tasks}

\ex{1} Calcula el ángulo que forman los planos:\begin{tasks}(2)
  \task $x+y=3$ \ y \ $x+z=1$
  \task $2x-y+2z=0$ \ y el plano $OXY$
\end{tasks}

\ex{2} Calcula los ángulos que forma la recta de vector director $(2,\,-1,\,2)$ con cada uno de los ejes de coordenadas.

\ex{2} Halla $m$ para que la recta $(x,\,y,\,z)=(1,\,0,\,0)+\lambda\,(1,\,m,\,1)$ forme un ángulo de $30^\circ$ con el plano $z=0$.

\ex{2} ¿Hay algún valor de $k$ para el que los planos $x+ky+z=0$ y $2x-y+kz=3$ sean perpendiculares? ¿Y paralelos?

\ex{2} \pau{Canarias, extraordinaria 2026} Se sabe que el plano $\pi:x+3y+2z=10$ y la recta \mbox{$r:\dfrac{x+21}{11}=3-y=\dfrac{5-z}{2}$} son secantes.\begin{tasks}(1)
  \task Halla el punto de corte del plano $\pi$ y la recta $r$.
  \task Calcula el ángulo que forman el plano $\pi$ y la recta $r$.
\end{tasks}

\ex{2} \pau{Castilla-La Mancha, ordinaria 2024} Sean los vectores $\vec u=(1,\,a,\,a)$ y $\vec v=(-1,\,0,\,2)$, con $a\in\mathbb R$. Determina el valor de $a$ para que el ángulo entre los vectores $\vec u$ y $\vec v$ sea de $60^\circ$.

\ex{2} \pau{País Vasco, extraordinaria 2025} Sean $\pi$ el plano que pasa por los puntos $A(2,\,3,\,4)$, $B(3,\,-1,\,-2)$, $C(5,\,-1,\,2)$ y $r$ la recta que pasa por los puntos $D(6,\,-5,\,-4)$, $E(7,\,1,\,4)$.\begin{tasks}(1)
  \task Calcula el ángulo entre el plano $\pi$ y la recta $r$, expresando el resultado en grados, minutos y segundos.
  \task En caso de que $r$ y $\pi$ se corten, calcula el punto de intersección. En caso contrario, calcula la distancia entre la recta $r$ y el plano $\pi$.
\end{tasks}

\ex{2} \pau{Canarias, extraordinaria 2024}\begin{tasks}(1)
  \task Justifica si pueden existir vectores $\vec u$ y $\vec v$, con el mismo origen, que cumplan $|\vec u|=2$, $|\vec v|=3$ y $\vec u\cdot\vec v=8$.
  \task Dados el plano $\pi:x+3y+2z+3=0$ y la recta $r:\left\{\begin{aligned}2x-3y-z&=4\\x+y+2z&=-3\end{aligned}\right.$, que son secantes, calcula el punto de corte de la recta y el plano y el ángulo que forman.
\end{tasks}

\ex{2} \pau{La Rioja, ordinaria 2024} Halla la ecuación continua de la recta $s$ que está contenida en el plano $\pi:x+y-2z+1=0$ y corta perpendicularmente a la recta $r:\left\{\begin{aligned}x+y+z&=-1\\4x-y+z&=3\end{aligned}\right.$


% =====================================================
\seccio{2}{Distancia entre puntos y de un punto a un plano}

\ex{1} Calcula:\begin{tasks}(1)
  \task la distancia del punto $P(2,\,-1,\,3)$ al plano $3x-6y+2z+1=0$;
  \task la distancia del origen al plano \mbox{$\dfrac x2+\dfrac y3+\dfrac z6=1$};
  \task la distancia entre los puntos $A(1,\,2,\,-1)$ y $B(3,\,-1,\,5)$.
\end{tasks}\recorda{\mbox{$d(P,\pi)=\dfrac{|Ax_0+By_0+Cz_0+D|}{\sqrt{A^2+B^2+C^2}}$}}

\ex{1} Halla la proyección ortogonal del punto $P(3,\,3,\,-1)$ sobre el plano $\pi:x+y-z=1$ y comprueba que la distancia de $P$ a la proyección es $d(P,\pi)$.

\ex{2} Halla los puntos del eje $OZ$ que distan 3 unidades del plano $2x+y-2z=1$.

\ex{2} Halla $a$ para que el punto $P(1,\,a,\,1)$ diste 2 unidades del plano $x+2y+2z-3=0$.

\ex{2} \pau{Asturias, ordinaria 2026} Se considera el plano $\pi:x+2y+3z=6$ y el punto $P(0,\,-1,\,-2)$.\begin{tasks}(1)
  \task Halla la ecuación de una recta paralela a $\pi$ que contenga el punto $P$. ¿Es única esta recta?
  \task Calcula el punto $Q\in\pi$ que dista menos de $P$. ¿Qué distancia hay entre $P$ y $Q$?
\end{tasks}

\ex{2} \pau{Cataluña, extraordinaria 2024} Considera el punto $P(1,\,3,\,0)$ y el plano $\pi:x+2y-2z=-7$.\begin{tasks}(1)
  \task Sea $r$ la recta perpendicular a $\pi$ que pasa por $P$. Calcula el punto de intersección de $\pi$ con $r$.
  \task Calcula la distancia $d$ del punto $P$ al plano $\pi$.
  \task Calcula la ecuación de otro plano $\pi'$ paralelo a $\pi$ que también esté a distancia $d$ de $P$.
\end{tasks}

\ex{2} \pau{Galicia, ordinaria 2026} Considera los puntos $P(1,\,0,\,1)$, $Q(0,\,1,\,0)$ y $R(1,\,1,\,1)$.\begin{tasks}(1)
  \task Demuestra que $P$, $Q$ y $R$ no están alineados y obtén la ecuación implícita del plano $\pi$ que pasa por ellos.
  \task Calcula el valor o valores de $a$ que hacen que la distancia del punto $S(a,\,0,\,0)$ a este plano sea igual a $\sqrt2$ unidades de longitud.
\end{tasks}

\ex{2} \pau{Asturias, extraordinaria 2026} Dado el plano $\pi:2x+2y+z=5$:\begin{tasks}(1)
  \task Halla la ecuación de un plano paralelo a $\pi$ cuya distancia al origen sea 2. ¿Hay más de un plano que cumpla estas condiciones?
  \task Calcula el punto $P$ del plano $\pi$ que está más próximo al origen $O$.
  \task Calcula el área del paralelogramo formado por los segmentos $OP$ y $OQ$, donde $Q(3,\,3,\,3)$.
\end{tasks}

\ex{2} \pau{Andalucía, ordinaria 2025} Sean los puntos $O(0,\,0,\,0)$, $A(0,\,2,\,-2)$, $B(1,\,2,\,m)$ y $C(2,\,3,\,2)$.\begin{tasks}(1)
  \task Halla los valores de $m$ para que el tetraedro determinado por los puntos $O$, $A$, $B$ y $C$ tenga un volumen de 3 unidades cúbicas.
  \task Para $m=0$, calcula la distancia del punto $O$ al plano que pasa por los puntos $A$, $B$ y $C$.
\end{tasks}


% =====================================================
\seccio{3}{Distancia de un punto a una recta}

\ex{1} Calcula la distancia:\begin{tasks}(1)
  \task del punto $P(1,\,2,\,3)$ a la recta $x=y=z$;
  \task del punto $P(0,\,1,\,-1)$ a la recta $r:\left\{\begin{aligned}x-y&=0\\z&=2\end{aligned}\right.$
\end{tasks}\recorda{\mbox{$d(P,r)=\dfrac{|\overrightarrow{AP}\times\vec d|}{|\vec d|}$}, con $A\in r$ y $\vec d$ un vector director de $r$.}

\ex{2} Halla el punto $Q$ de la recta $r:(x,\,y,\,z)=(1,\,0,\,1)+\lambda\,(1,\,1,\,1)$ más próximo a $P(3,\,1,\,4)$, la distancia de $P$ a $r$ y la recta que pasa por $P$ y corta perpendicularmente a $r$.

\ex{2} Halla los puntos de la recta $r:(x,\,y,\,z)=(1,\,1,\,0)+\lambda\,(2,\,1,\,2)$ que están a 3 unidades del punto $A(1,\,1,\,0)$.

\ex{2} Calcula la altura del triángulo de vértices $A(1,\,0,\,0)$, $B(0,\,1,\,0)$ y $C(1,\,1,\,1)$ correspondiente al lado $AB$.

\ex{2} \pau{La Rioja, extraordinaria 2026} Sean la recta \mbox{$r:\dfrac x3=y=\dfrac{z-11}{-1}$} y el punto $P(0,\,1,\,1)$.\begin{tasks}(1)
  \task Determina la ecuación del plano que pasa por el punto $P$ y es perpendicular a la recta $r$.
  \task Halla la distancia de $P$ a la recta $r$.
\end{tasks}

\ex{2} \pau{Cantabria, ordinaria 2025} Considera la recta $r:\left\{\begin{aligned}2x-y&=3\\y-2z&=1\end{aligned}\right.$ \ y el punto $P(1,\,1,\,1)$.\begin{tasks}(1)
  \task Determina los puntos de $r$ que están a una distancia de $\sqrt{14}$ unidades de $P$.
  \task Obtén la ecuación del plano que contiene $r$ y $P$.
  \task Calcula la distancia entre $r$ y $P$.
\end{tasks}

\ex{2} \pau{Extremadura, extraordinaria 2025} Dados la recta \mbox{$s:\dfrac{x}{-4}=\dfrac{y-1}{3}=\dfrac{z+1}{0}$}, el plano $\pi:x-2y+3z-6=0$ y el punto $P(1,\,-1,\,0)$:\begin{tasks}(1)
  \task Obtén la ecuación del plano perpendicular a la recta $s$ que pasa por $P$.
  \task Calcula la distancia del punto $P$ a la recta $s$.
  \task Calcula el ángulo que forma la recta $s$ con el plano $\pi$.
\end{tasks}

\ex{2} \pau{Cantabria, ordinaria 2026} Considera los puntos $A(0,\,0,\,4)$, $B(-3,\,0,\,0)$ y $C(0,\,-2,\,0)$.\begin{tasks}(1)
  \task Halla la ecuación general del plano que contiene $A$, $B$ y $C$.
  \task Calcula la distancia entre el punto $P(1,\,1,\,1)$ y la recta que contiene $A$ y $B$.
  \task Calcula el volumen del tetraedro de vértices $A$, $B$, $C$ y $D(0,\,0,\,0)$.
\end{tasks}


% =====================================================
\seccio{4}{Distancia entre dos rectas}

\ex{1} Comprueba que las rectas $r:(x,\,y,\,z)=(1,\,1,\,0)+\lambda\,(1,\,-1,\,1)$ y $s:(x,\,y,\,z)=(2,\,0,\,3)+\mu\,(-1,\,1,\,-1)$ son paralelas y calcula la distancia entre ellas.

\ex{2} \label{ex:pc}Comprueba que las rectas $r:\left\{\begin{aligned}x-y&=0\\z&=0\end{aligned}\right.$ \ y \ $s:\left\{\begin{aligned}x&=1\\y-z&=0\end{aligned}\right.$ \ se cruzan y calcula la distancia entre ellas.\recorda{Si se cruzan, \mbox{$d(r,s)=\dfrac{\big|[\overrightarrow{PQ},\vec d_r,\vec d_s]\big|}{|\vec d_r\times\vec d_s|}$}.}

\ex{2} Halla $a$ para que la distancia entre el eje $OX$ y la recta $s:(x,\,y,\,z)=(0,\,1,\,a)+\mu\,(1,\,1,\,0)$ sea 2.

\ex{2} \pau{País Vasco, ordinaria 2025} Se consideran la recta $r$ que pasa por los puntos $A(-4,\,4,\,1)$ y $B(1,\,0,\,-1)$ y la recta $s$ que pasa por los puntos $C(-2,\,7,\,-1)$ y $D(2,\,3,\,-1)$.\begin{tasks}(1)
  \task Calcula la posición relativa de la recta $r$ respecto de la recta $s$.
  \task En caso de que sean paralelas o se crucen, calcula la distancia entre ellas. Si se cortan, calcula el punto de intersección.
\end{tasks}

\ex{2} \pau{La Rioja, extraordinaria 2025} Calcula la distancia entre las rectas $r:\left\{\begin{aligned}x&=5+\lambda\\y&=-1\\z&=8+2\lambda\end{aligned}\right.$ \ y \ $s:\left\{\begin{aligned}x&=2+3\mu\\y&=2-\mu\\z&=-1+4\mu\end{aligned}\right.$


% =====================================================
\seccio{5}{Distancia entre recta y plano y entre dos planos}

\ex{1} Comprueba que la recta $r:x=2+3\lambda,\ y=-\lambda,\ z=1+\lambda$ es paralela al plano $\pi:x+2y-z+4=0$ y calcula la distancia entre ellos.

\ex{1} Calcula la distancia entre los planos $3x-4y+12z=5$ y $6x-8y+24z+16=0$.

\ex{2} Halla los planos paralelos a $2x+y-2z+1=0$ que distan 3 unidades del origen.

\ex{2} Halla $m$ para que la recta $(x,\,y,\,z)=(1,\,1,\,1)+\lambda\,(m,\,1,\,2)$ sea paralela al plano $x+2y-2z=0$ y calcula entonces la distancia entre la recta y el plano.

\ex{2} \pau{Andalucía, extraordinaria 2025} Considera la recta $r:\left\{\begin{aligned}x-y+z&=3\\x+2y-z&=4\end{aligned}\right.$ \ y el plano $\pi:mx-y-2z=5$.\begin{tasks}(1)
  \task Halla $m$ para que $r$ y $\pi$ sean paralelos.
  \task Para $m=-8$, calcula la distancia de la recta $r$ al plano $\pi$.
\end{tasks}

\ex{2} \pau{Extremadura, ordinaria 2025}\begin{tasks}(1)
  \task Comprueba que el plano $\pi:x+y-z=3$ y la recta \mbox{$r:\dfrac x3=\dfrac{y-1}{-1}=\dfrac{z+1}{2}$} no se cortan.
  \task Calcula la distancia entre el plano $\pi$ y la recta $r$.
  \task Obtén la ecuación del plano perpendicular a la recta $r$ que pasa por el punto $(0,\,1,\,-1)$.
\end{tasks}

\ex{2} \pau{Galicia, extraordinaria 2026} Considera la recta \mbox{$r:\dfrac{x-1}{2}=\dfrac{y-1}{1}=\dfrac z3$}.\begin{tasks}(1)
  \task Calcula la ecuación implícita o general del plano $\pi$ que contiene la recta $r$ y el punto $Q(2,\,-1,\,1)$.
  \task Comprueba que la recta $r$ es paralela al plano $\pi':x+y-z-1=0$. Calcula la distancia de $r$ a $\pi'$.
\end{tasks}

\ex{2} \pau{Murcia, ordinaria 2026} Considera el plano $\pi:2x-y+3z-5=0$ y la recta \mbox{$r:\dfrac{x-1}{2}=\dfrac{y+1}{1}=\dfrac{z-2}{-1}$}.\begin{tasks}(1)
  \task Estudia la posición relativa de $r$ y $\pi$.
  \task Si $r$ corta a $\pi$ en un solo punto, calcula este punto y el ángulo que forma $r$ con $\pi$. Si $r$ es paralela a $\pi$, calcula la distancia de $r$ a $\pi$.
\end{tasks}

\ex{2} \pau{Castilla y León, extraordinaria 2025} Sean el plano $\pi:x+y-z=2$ y la recta \mbox{$r:\dfrac{x-1}{-3}=\dfrac y2=\dfrac{z-2}{-1}$}. Calcula la ecuación de un plano $\pi'$ paralelo al plano $\pi$ que esté a una distancia de $2\sqrt3$ unidades de la recta $r$. ¿Es único este plano? Justifica la respuesta.

\ex{2} \pau{Cataluña, extraordinaria 2025} Considera el plano $\pi:2x-y+z=5$ y el punto $P(0,\,1,\,3)$.\begin{tasks}(1)
  \task Comprueba que la distancia del punto $P$ al plano $\pi$ es \mbox{$\dfrac{\sqrt6}{2}$}.
  \task Halla la ecuación general de un plano $\pi_1$ paralelo a $\pi$ que pase por el punto $P$. ¿Cuánto vale la distancia entre $\pi_1$ y $\pi$?
  \task Halla la ecuación general de un segundo plano $\pi_2$, distinto de $\pi_1$, que sea paralelo a $\pi$ y que esté a una distancia \mbox{$\dfrac{\sqrt6}{2}$} de $\pi$.
\end{tasks}


% =====================================================
\seccio{6}{Recta que pasa por un punto y corta a dos rectas}

\ex{1} Halla la recta que pasa por el punto $P(2,\,1,\,1)$ y corta a las rectas $r:\left\{\begin{aligned}x&=0\\y&=0\end{aligned}\right.$ \ y \ $s:\left\{\begin{aligned}x&=1\\z&=0\end{aligned}\right.$ \ Calcula también los puntos de corte.\recorda{$t=\pi_1\cap\pi_2$, donde $\pi_1$ es el plano que contiene $P$ y $r$, y $\pi_2$, el que contiene $P$ y $s$. Comprueba que $t$ no es paralela a $r$ ni a $s$.}

\ex{2} Halla la recta paralela al vector $\vec v=(1,\,0,\,1)$ que corta a las rectas $r$ y $s$ del ejercicio anterior.

\ex{2} Halla la recta contenida en el plano $\pi:x+y+z=3$ que pasa por el punto $P(1,\,1,\,1)$ y corta a la recta $r:(x,\,y,\,z)=\lambda\,(1,\,2,\,0)$.

\ex{2} Razona que no existe ninguna recta que pase por el origen y corte a las rectas $r:(x,\,y,\,z)=(1,\,0,\,0)+\lambda\,(0,\,1,\,0)$ y $s:(x,\,y,\,z)=(0,\,0,\,1)+\mu\,(1,\,0,\,0)$.

\ex{2} \pau{Navarra, extraordinaria 2026} Halla la ecuación continua de la recta $t$ que pasa por el punto $A(2,\,4,\,-2)$ y corta a las rectas \mbox{$s:\dfrac{x-2}{-1}=\dfrac{y+1}{2}=\dfrac z1$} \ y \ $r:\left\{\begin{aligned}x&=-1+\lambda\\y&=-2\\z&=1+3\lambda\end{aligned}\right.$

\ex{2} \pau{Castilla y León, extraordinaria 2026} Dados el plano $\pi:x+y=0$, el punto $A(1,\,-1,\,0)$ y la recta $s:x=y=z$, halla la recta $r$, en forma paramétrica, que pasa por $A$, está contenida en $\pi$ y corta a la recta $s$.

\ex{2} \pau{Madrid, extraordinaria 2024} Dados el punto $P(5,\,-1,\,2)$ y las rectas \mbox{$r:\dfrac{x-2}{3}=\dfrac{y+1}{-1}=\dfrac{z}{1}$} \ y \ $s:\left\{\begin{aligned}x-y&=5\\x+z&=3\end{aligned}\right.$\begin{tasks}(1)
  \task Estudia la posición relativa de las dos rectas y halla la distancia entre ellas.
  \task Determina una ecuación de la recta que pasa por $P$ y corta perpendicularmente a la recta $r$.
\end{tasks}


% =====================================================
\seccio{7}{Perpendicular común}

\ex{2} Halla los puntos $X\in r$ e $Y\in s$ más próximos de las rectas $r:\left\{\begin{aligned}x-y&=0\\z&=0\end{aligned}\right.$ \ y \ $s:\left\{\begin{aligned}x&=1\\y-z&=0\end{aligned}\right.$ \ del ejercicio \ref{ex:pc} y su perpendicular común. Comprueba que $|\overrightarrow{XY}|=d(r,s)$.\recorda{Con los puntos genéricos $X\in r$, $Y\in s$, impón $\overrightarrow{XY}\perp\vec d_r$ y $\overrightarrow{XY}\perp\vec d_s$.}

\ex{2} Sean $r:(x,\,y,\,z)=(1,\,0,\,0)+\lambda\,(1,\,1,\,0)$ y $s:(x,\,y,\,z)=(0,\,0,\,1)+\mu\,(1,\,-1,\,0)$. Halla la perpendicular común por el método de los planos y comprueba el resultado con los puntos genéricos. ¿Cuál es la distancia entre $r$ y $s$?\recorda{$\vec w=\vec d_r\times\vec d_s$; \ $\pi_r$ contiene $r$ y $\vec w$; \ $\pi_s$ contiene $s$ y $\vec w$; \ la perpendicular común es $\pi_r\cap\pi_s$.}

\ex{2} \pau{Murcia, ordinaria 2025} Considera las rectas $r:\left\{\begin{aligned}x&=1\\y&=5+\lambda\\z&=\lambda\end{aligned}\right.$ \ y \ $s:\left\{\begin{aligned}x&=7+3\mu\\y&=-5+\mu\\z&=7\end{aligned}\right.$\begin{tasks}(1)
  \task Comprueba que se cruzan.
  \task Halla la ecuación de la recta perpendicular a ambas.
  \task Calcula la distancia entre $r$ y $s$.
\end{tasks}

\ex{2} \pau{Navarra, ordinaria 2026} Halla la ecuación continua de la recta $t$ que corta perpendicularmente a las rectas \mbox{$r:\dfrac x1=\dfrac{y+3}{1}=\dfrac{z-2}{-2}$} \ y \ $s:\left\{\begin{aligned}x&=-4+3\lambda\\y&=1+\lambda\\z&=-1+2\lambda\end{aligned}\right.$

\ex{2} \pau{Navarra, extraordinaria 2025} Calcula la ecuación continua de la recta $t$ que corta perpendicularmente a las rectas $r:\left\{\begin{aligned}x+y-2z-10&=0\\x+z-3&=0\end{aligned}\right.$ \ y \ \mbox{$s:\dfrac{x-4}{-1}=\dfrac{y+1}{1}=\dfrac{z}{-1}$}


% =====================================================
\seccio{8}{Simétricos respecto de un punto y de una recta}

\ex{1} Calcula el simétrico de $A(1,\,-2,\,3)$ respecto del punto $B(0,\,1,\,-1)$. Halla también el punto $M$ respecto del cual $A$ y $A'(5,\,0,\,1)$ son simétricos.

\ex{2} Calcula el simétrico del punto $P(2,\,1,\,0)$ respecto de la recta $r:(x,\,y,\,z)=(0,\,0,\,1)+\lambda\,(1,\,0,\,1)$.\recorda{Plano por $P$ perpendicular a $r$ $\to$ $M=r\cap\pi$ $\to$ $P'=2M-P$.}

\ex{2} Calcula el simétrico de:\begin{tasks}(1)
  \task $P(1,\,2,\,3)$ respecto del eje $OX$;
  \task $P(0,\,0,\,1)$ respecto de la recta $r:\left\{\begin{aligned}x+y&=2\\z&=0\end{aligned}\right.$
\end{tasks}

\ex{2} \pau{Castilla y León, ordinaria 2024} Halla el punto simétrico del punto $P(1,\,0,\,-1)$ respecto de la recta \mbox{$r:\dfrac{x-1}{1}=\dfrac y2=\dfrac z2$}.

\ex{2} \pau{País Vasco, extraordinaria 2025} Halla el punto $P'$ simétrico de $P(4,\,-3,\,0)$ respecto de la recta $r$ que pasa por los puntos $A(-2,\,0,\,1)$ y $B(-2,\,-1,\,0)$.


% =====================================================
\seccio{9}{Simétricos respecto de un plano}

\ex{1} Calcula el simétrico del punto $P(3,\,1,\,2)$ respecto del plano $\pi:2x-y+z-1=0$.\recorda{Recta por $P$ con vector director $\vec n$ $\to$ $M=s\cap\pi$ (la proyección) $\to$ $P'=2M-P$.}

\ex{1} Calcula los simétricos del punto $A(1,\,2,\,3)$ respecto de los planos $z=0$ y $x=y$.

\ex{2} Halla el plano respecto del cual son simétricos los puntos $P(1,\,2,\,-1)$ y $P'(3,\,0,\,1)$.

\ex{3} Sean $r:(x,\,y,\,z)=(1,\,1,\,1)+\lambda\,(1,\,2,\,0)$ y $\pi:x+y+z=0$.\begin{tasks}(1)
  \task Calcula la recta simétrica de $r$ respecto de $\pi$.
  \task Calcula la proyección ortogonal de $r$ sobre $\pi$.
\end{tasks}

\ex{2} Calcula la recta simétrica de $r:(x,\,y,\,z)=(2,\,0,\,0)+\lambda\,(0,\,1,\,-1)$ respecto del plano $y+z=4$.

\ex{2} \pau{Andalucía, ordinaria 2024}\begin{tasks}(1)
  \task Halla el punto simétrico de $P(2,\,2,\,1)$ respecto de la recta $r:\left\{\begin{aligned}x-2y+z&=2\\y-z&=1\end{aligned}\right.$
  \task Halla el punto simétrico de $Q(1,\,-1,\,-3)$ respecto del plano $\pi:x-2y+z+6=0$.
\end{tasks}

\ex{2} \pau{Andalucía, extraordinaria 2026} Considera la recta $r:\left\{\begin{aligned}x&=-2\lambda\\y&=-1+2\lambda\\z&=3+2\lambda\end{aligned}\right.$ \ y el plano $\pi:x-2y+z-3=0$.\begin{tasks}(1)
  \task Determina la ecuación de la recta que corta perpendicularmente a $r$ y está contenida en $\pi$.
  \task Halla el punto simétrico de $P(1,\,1,\,1)$ respecto de $\pi$.
\end{tasks}

\ex{2} \pau{Canarias, ordinaria 2025} Se consideran el punto $A(1,\,0,\,2)$, el plano $\pi:-x+2y+z+1=0$ y la recta $r:\left\{\begin{aligned}4x-7y+2z&=7\\y-2z+5&=0\end{aligned}\right.$\begin{tasks}(1)
  \task Halla la posición relativa del plano $\pi$ y la recta $r$.
  \task Halla el punto simétrico de $A$ respecto del plano $\pi$.
\end{tasks}

\ex{2} \pau{Extremadura, extraordinaria 2026} Dados los planos $\pi_1:x+ky-3z-6=0$, $\pi_2:2x-y-5=0$ y el punto $A(2,\,3,\,-4)$:\begin{tasks}(1)
  \task Calcula el valor de $k$ para que los planos $\pi_1$ y $\pi_2$ sean perpendiculares.
  \task Para $k=2$, calcula el simétrico del punto $A$ respecto del plano $\pi_1$.
  \task Para $k=2$, calcula la distancia del punto $A$ al plano $\pi_1$.
\end{tasks}

\ex{2} \pau{País Vasco, extraordinaria 2026} Sea $P$ el punto de corte de las rectas $r_1:\left\{\begin{aligned}x&=1+t\\y&=0\\z&=-t\end{aligned}\right.$ \ y \ $r_2:\left\{\begin{aligned}x&=s\\y&=-1+s\\z&=0\end{aligned}\right.$ \ Calcula el punto $P'$ simétrico de $P$ respecto del plano de ecuación $2x+y-z=3$.

\ex{2} \pau{Madrid, ordinaria 2026} Dados el plano $\pi:2x+2y-z=13$ y la recta \mbox{$r:\dfrac{x-2}{1}=\dfrac y1=\dfrac z4$}:\begin{tasks}(1)
  \task Halla el punto simétrico del punto $P(2,\,0,\,0)$ respecto del plano $\pi$.
  \task Halla la distancia entre el plano $\pi$ y la recta $r$.
\end{tasks}

\ex{2} \pau{Galicia, extraordinaria 2026} Considera el punto $P(3,\,0,\,-1)$.\begin{tasks}(1)
  \task Calcula el punto $P'$ simétrico de $P$ respecto del punto $Q(2,\,-1,\,2)$.
  \task Calcula el punto $P''$ simétrico de $P$ respecto del plano $\pi:x+2y+z-8=0$.
\end{tasks}


% =====================================================
\seccio{10}{Áreas y volúmenes}

\ex{1} Considera el triángulo de vértices $A(1,\,2,\,0)$, $B(3,\,0,\,1)$ y $C(0,\,1,\,2)$.\begin{tasks}(1)
  \task Calcula su área.
  \task Calcula el ángulo del vértice $A$.
  \task Calcula la altura correspondiente al lado $AB$.
\end{tasks}\recorda{$S=\tfrac12|\overrightarrow{AB}\times\overrightarrow{AC}|$. Para el ángulo de un triángulo no ponemos valor absoluto al producto escalar.}

\ex{2} Calcula el volumen del tetraedro de vértices $A(1,\,0,\,0)$, $B(0,\,1,\,0)$, $C(0,\,0,\,1)$ y $D(2,\,2,\,2)$, y su altura desde $D$.

\ex{2} Halla $a$ para que el tetraedro de vértices $O$, $(1,\,0,\,0)$, $(0,\,1,\,0)$ y $(1,\,1,\,a)$ tenga volumen 2.

\ex{2} Tres vértices consecutivos de un paralelogramo $ABCD$ son $A(1,\,0,\,1)$, $B(2,\,1,\,0)$ y $C(4,\,2,\,1)$. Halla el cuarto vértice y el área del paralelogramo.

\ex{2} El plano $\pi:2x+3y+6z=12$ corta los ejes en $A$, $B$ y $C$. Calcula el volumen del tetraedro $OABC$, la distancia del origen a $\pi$ y, con estos dos datos, el área del triángulo $ABC$.

\ex{2} \pau{Cantabria, extraordinaria 2026} Considera los puntos $A(0,\,5,\,-1)$ y $B(-3,\,11,\,-10)$ y la recta \mbox{$r:\dfrac{x-3}{1}=\dfrac{y+2}{-2}=\dfrac z3$}.\begin{tasks}(1)
  \task Estudia la posición relativa de $r$ y la recta que contiene $A$ y $B$.
  \task Calcula el área del triángulo formado por los vértices $A$, $B$ y un punto genérico de $r$.
  \task Halla la ecuación general o implícita del plano que contiene $r$ y $A$.
\end{tasks}

\ex{2} \pau{Cantabria, ordinaria 2024} Sean $A(6,\,2,\,-1)$, $B(3,\,0,\,5)$ y $C(-2,\,1,\,2)$ los vértices de un triángulo.\begin{tasks}(1)
  \task Calcula los ángulos internos del triángulo.
  \task Calcula el área del triángulo.
\end{tasks}

\ex{2} \pau{Extremadura, ordinaria 2024} Dados los puntos $A(1,\,2,\,1)$, $B(0,\,3,\,1)$ y $C(1,\,0,\,-1)$, determina:\begin{tasks}(1)
  \task un vector unitario y ortogonal a los vectores $\overrightarrow{AB}$ y $\overrightarrow{AC}$;
  \task el ángulo determinado por estos vectores;
  \task el área del triángulo que forman $A$, $B$ y $C$.
\end{tasks}

\ex{2} \pau{Madrid, ordinaria 2025} Sean el punto $P(0,\,1,\,1)$ y el plano $\pi:x+y=2$.\begin{tasks}(1)
  \task Halla la distancia del punto $P$ al plano $\pi$.
  \task Determina el punto $Q$ del plano $\pi$ cuya distancia a $P$ es igual a la distancia de $P$ a $\pi$.
  \task Halla el área del triángulo formado por $P$ y los puntos de corte del plano $\pi$ con los ejes de coordenadas.
\end{tasks}

\ex{2} \pau{Navarra, extraordinaria 2024} El punto $P(4,\,5,\,0)$ es el punto medio de un lado de un cuadrado. El lado paralelo al anterior está contenido en la recta $r:\left\{\begin{aligned}2x+2y+z&=0\\2x-z-2&=0\end{aligned}\right.$ \ Calcula los dos vértices que determinan este segundo lado.


% =====================================================
\seccio{11}{Lugares geométricos}

\ex{1} Halla el plano mediador del segmento de extremos $A(2,\,1,\,-1)$ y $B(0,\,3,\,1)$ de dos maneras: imponiendo $d(X,A)=d(X,B)$ y con el punto medio y el vector $\overrightarrow{AB}$.

\ex{2} Halla los planos bisectores de los planos $\pi_1:x+y-z=1$ y $\pi_2:x-y+z=3$.\recorda{Son los puntos que equidistan de los dos planos: \mbox{$\dfrac{|A_1x+B_1y+C_1z+D_1|}{|\vec n_1|}=\dfrac{|A_2x+B_2y+C_2z+D_2|}{|\vec n_2|}$}.}

\ex{1} Halla el lugar geométrico de los puntos que equidistan de los planos $2x-y+2z=1$ y $2x-y+2z=7$.

\ex{2} Escribe la ecuación de la esfera de centro $C(1,\,0,\,-2)$ que pasa por el punto $P(3,\,2,\,-1)$ y la del plano tangente a la esfera en $P$.

\ex{2} Halla la esfera de centro $C(2,\,-1,\,1)$ tangente al plano $x+2y+2z+7=0$ y el punto de tangencia.

\ex{3} Halla el lugar geométrico de los puntos del espacio que equidistan de $A(1,\,0,\,0)$, $B(0,\,1,\,0)$ y $C(0,\,0,\,1)$.

\ex{2} \pau{Andalucía, extraordinaria 2024} Considera los puntos $A(4,\,0,\,0)$ y $B(0,\,2,\,0)$. Calcula los puntos del plano $OXZ$ que forman un triángulo equilátero con $A$ y $B$.

\ex{2} \pau{Andalucía, ordinaria 2024} Considera las rectas $r:\left\{\begin{aligned}y&=0\\2x-z&=0\end{aligned}\right.$ \ y \ $s:\left\{\begin{aligned}x+y+7&=0\\z&=0\end{aligned}\right.$\begin{tasks}(1)
  \task Estudia la posición relativa de $r$ y $s$.
  \task Calcula la ecuación del plano paralelo a $r$ y a $s$ que equidista de ambas rectas.
\end{tasks}

\ex{2} \pau{Castilla y León, ordinaria 2026} Dados el plano $\pi_1:2x-3y+z=a$ y el plano $\pi_2$ determinado por el punto $P(0,\,2,\,4)$ y los vectores directores $\vec v_1=(0,\,2,\,6)$ y $\vec v_2=(1,\,0,\,b)$:\begin{tasks}(1)
  \task Calcula los valores de $a$ y $b$ para que $\pi_1$ y $\pi_2$ sean paralelos y distintos.
  \task Para $a=1$ y $b=0$, determina las ecuaciones paramétricas de la recta intersección de $\pi_1$ y $\pi_2$.
  \task Para $a=4$ y $b=-2$, determina los puntos que están a igual distancia de $\pi_1$ y $\pi_2$.
\end{tasks}

\ex{2} \pau{Cataluña, ordinaria 2025} Considera el plano $\pi:x+y=0$.\begin{tasks}(1)
  \task Calcula la ecuación del plano $\pi'$ que es perpendicular a $\pi$ y contiene los puntos $P(1,\,-1,\,2)$ y $Q(3,\,-3,\,6)$.
  \task Calcula la ecuación paramétrica de la recta contenida en $\pi'$ que contiene los puntos de $\pi'$ que están a la misma distancia de $P$ que de $Q$.
\end{tasks}

\ex{2} \pau{Cataluña, ordinaria 2024} Considera los puntos $A(1,\,2,\,3)$ y $B(-3,\,-2,\,3)$.\begin{tasks}(1)
  \task Calcula la ecuación del plano $\pi$ que es perpendicular a la recta $AB$ y que pasa por el punto medio entre $A$ y $B$. Justifica que este plano está formado, precisamente, por los puntos $P$ que están a igual distancia de $A$ que de $B$.
  \task Calcula las distancias de $A$ y de $B$ al plano $\pi$ y comprueba que son iguales. ¿Es una casualidad?
  \task Sea $C(-7,\,6,\,3)$. ¿El triángulo $ABC$ es isósceles? Calcula su área.
\end{tasks}

\ex{2} \pau{Murcia, extraordinaria 2024} Considera los planos $x+y+z=-3$ y $x+y-z=3$ y la recta \mbox{$r:\dfrac{x-1}{2}=\dfrac{y+1}{1}=\dfrac z3$}.\begin{tasks}(1)
  \task Comprueba que ambos planos se cortan y calcula el ángulo que forman.
  \task Estudia la posición relativa de la recta $r$ con el plano $x+y-z=3$.
  \task Determina los puntos de la recta $r$ que equidistan de ambos planos.
\end{tasks}

\ex{2} \pau{Navarra, extraordinaria 2025} Calcula la ecuación del plano que equidista de los puntos $A(3,\,3,\,5)$ y $B(1,\,-1,\,1)$. Calcula la ecuación de los planos que distan 6 unidades del plano anterior.

\ex{2} \pau{Madrid, extraordinaria 2024} Sean los puntos $P(1,\,-1,\,3)$ y $Q(2,\,1,\,-1)$.\begin{tasks}(1)
  \task Determina una ecuación del plano respecto del cual ambos puntos son simétricos.
  \task El segmento $PQ$ es uno de los tres lados de un triángulo en el que la suma de los cuadrados de las longitudes de los lados es 34 y el tercer vértice se encuentra en la recta $r:x-2=y=z$. Calcula las coordenadas del tercer vértice sabiendo que ninguna de sus coordenadas es nula.
\end{tasks}

\ex{2} \pau{País Vasco, ordinaria 2026} Sean los puntos $A(1,\,1,\,2)$ y $B(-1,\,0,\,3)$.\begin{tasks}(1)
  \task Halla la ecuación del plano $\pi$ que cumple que $A$ y $B$ son simétricos respecto de $\pi$.
  \task ¿Cuál es la distancia de $A$ a $\pi$?
\end{tasks}


% =====================================================
\seccio{12}{Problemas y cuestiones}

\ex{2} Un rayo de luz sale del punto $P(1,\,2,\,5)$ en la dirección del vector $(1,\,0,\,-1)$, se refleja en un espejo situado en el plano $z=0$ y continúa. Halla el punto donde toca el espejo y la ecuación del rayo reflejado.\recorda{El rayo reflejado es la recta que pasa por el punto de incidencia y por el simétrico de $P$ respecto del espejo.}

\ex{3} Sean $A(1,\,1,\,2)$ y $B(4,\,5,\,1)$, ambos por encima del plano $z=0$. Halla el punto $X$ del plano $z=0$ que hace mínima la suma $d(A,X)+d(X,B)$ y el valor de esta suma mínima.

\ex{2} Razona si son ciertas o falsas las afirmaciones siguientes:\begin{tasks}(1)
  \task Dos planos perpendiculares a un mismo plano son siempre paralelos entre sí.
  \task Si una recta es perpendicular a un plano, es perpendicular a todas las rectas contenidas en el plano.
  \task La distancia entre dos rectas que se cruzan puede ser 0.
  \task Si $d(P,\pi)=0$, el punto $P$ está en el plano $\pi$.
\end{tasks}

\ex{2} \pau{Asturias, ordinaria 2026} Un móvil se desplaza en línea recta desde el punto $P(2,\,1,\,3)$ hacia el punto $Q(0,\,-1,\,1)$. Cuando llega a la mitad de la distancia de $P$ a $Q$ hace un giro y se desplaza por la recta que pasa por $H(1,\,1,\,1)$.\begin{tasks}(1)
  \task Calcula las coordenadas del punto en el que hace el giro.
  \task Calcula la recta por la que se desplaza el móvil después del giro.
  \task Calcula el ángulo de giro del móvil.
\end{tasks}

\ex{2} \pau{Murcia, extraordinaria 2025} Un helicóptero situado en el punto $P(1,\,2,\,1)$ quiere aterrizar en el plano $\pi:x+y+3z=0$.\begin{tasks}(1)
  \task Calcula la ecuación en forma continua de la recta de la trayectoria que lo lleve al punto más próximo de $\pi$.
  \task Halla este punto.
  \task Calcula la distancia que debe recorrer.
\end{tasks}

\ex{2} \pau{Asturias, ordinaria 2025} Se están construyendo dos puentes rectos en un tramo de autovía, uno para cada carril. Los puentes siguen las ecuaciones $r_1(t)=(2+t,\,-1-2t,\,3+2t)$ y $r_2(s)=(1+2s,\,4-s,\,4-2s)$.\begin{tasks}(1)
  \task Estudia si los puentes son paralelos, se cortan o se cruzan.
  \task La empresa quiere construir un puente de servicio que los una y que sea lo más corto posible. ¿Qué longitud tendrá la vía de servicio? Indica sus puntos inicial y final.
\end{tasks}

\ex{2} \pau{Madrid, extraordinaria 2026} Se consideran el punto $P(1,\,2,\,1)$ y el plano de ecuación $2x-2y+z=1$. Se traza un cubo con una cara en el plano y un vértice en el punto $P$.\begin{tasks}(1)
  \task Determina la longitud del lado del cubo y su volumen.
  \task Halla otro vértice del cubo.
\end{tasks}

\ex{3} \pau{Aragón, extraordinaria 2025} Se quiere construir una estructura en forma de tetraedro cuya base tiene como vértices los puntos $A(0,\,0,\,0)$, $B\left(2,\,0,\,\tfrac12\right)$ y $C\left(\tfrac32,\,3,\,1\right)$; el vértice superior, $D$, se encuentra en una viga recta entre los puntos $E(0,\,1,\,3)$ y $F(3,\,2,\,3)$, es decir, $D=E+\lambda\,\overrightarrow{EF}$ con $\lambda\in[0,1]$.\begin{tasks}(1)
  \task Calcula el volumen máximo de esta estructura (todos los datos están en metros).
  \task Teniendo en cuenta que el volumen de una pirámide es un tercio del área de la base por la altura, calcula la altura de la estructura (desde $D$ hasta la base) si tomamos $\lambda=1$.
\end{tasks}

\ex{2} \pau{Castilla-La Mancha, ordinaria 2026} En un parque público se está instalando un sistema de vigilancia mediante sensores de movimiento a ambos lados de un muro recto, representado por el plano $\pi:x+y+z-6=0$. Uno de los sensores está situado en el punto $P(1,\,2,\,6)$ y se quiere instalar un segundo sensor en el punto $P'$, simétrico del anterior respecto del plano $\pi$.\begin{tasks}(1)
  \task Determina razonadamente las coordenadas del punto de instalación del segundo sensor, $P'$.
  \task Calcula la distancia entre los dos sensores.
\end{tasks}

\ex{2} \pau{Castilla-La Mancha, extraordinaria 2025} Para las fiestas del Corpus de Toledo se instalan toldos en las calles por donde pasa la procesión. En uno de ellos, los operarios colocan los puntos de apoyo $A(0,\,1,\,-2)$, $B(1,\,2,\,0)$, $C(0,\,0,\,1)$ y $D(1,\,0,\,m)$, con $m\in\mathbb R$.\begin{tasks}(1)
  \task Calcula el valor de $m$ para que los cuatro puntos sean coplanarios.
  \task Determina la ecuación del plano $\pi$ que contiene el toldo.
  \task Si los adornos florales deben estar como mínimo a 1 metro de distancia del toldo y se ha colocado un adorno de flores en el punto $P(1,\,2,\,3)$, ¿estará correctamente situado?
\end{tasks}

\ex{2} \pau{Islas Baleares, ordinaria 2024} Consideramos los puntos $A(0,\,0,\,0)$, $B(2,\,-1,\,3)$ y $C(-1,\,2,\,1)$.\begin{tasks}(1)
  \task Calcula el punto $D$ tal que $ABDC$ sea un paralelogramo.
  \task Calcula uno de los puntos $E$ del espacio de manera que la recta $AE$ sea perpendicular al plano $ABC$ y la distancia entre los puntos $A$ y $E$ sea 1.
  \task Escribe la ecuación de uno de los planos paralelos al plano $ABC$ que distan una unidad de este.
\end{tasks}

\vfill
\begin{center}\color{muted}\small\faCheckCircle[regular]\ \ ¡Ánimo! Comprueba tus resultados con el laboratorio métrico de los apuntes interactivos.\end{center}

\end{document}
